% !TeX program = pdflatex
%==============================================================================
%  Lista Teórica 11 --- KNN e Árvores de Classificação
%  O gabarito sai de "Lista teorica 11 - gabarito.tex".
%==============================================================================
\providecommand{\gabaritoopt}{}
\documentclass[11pt,a4paper]{article}
\usepackage[\gabaritoopt]{../../recursos/latex/estilo-lista}

\begin{document}
\cabecalholista{11}{KNN e Árvores de Classificação}

%------------------------------------------------------------------------------
\begin{exercicio}[o corte que o erro de classificação não enxerga]
Um nó tem $400$ observações, das quais $300$ são da classe~1 --- logo
$p=0{,}75$. Um corte candidato o divide em dois filhos de $200$ observações
cada:
\begin{center}\small
\renewcommand{\arraystretch}{1.3}
\begin{tabular}{@{}lccc@{}}
\toprule
        & $n$ & classe 1 & classe 0 \\
\midrule
pai     & $400$ & $300$ & $100$ \\
esquerda& $200$ & $100$ & $100$ \\
direita & $200$ & $200$ & $0$ \\
\bottomrule
\end{tabular}
\end{center}
As três medidas de impureza de um nó com proporção $p$ são
\[
  \text{erro} = \min(p,\,1-p), \qquad
  \text{Gini} = 2p(1-p), \qquad
  \text{entropia} = -p\log_2 p - (1-p)\log_2(1-p).
\]
\begin{enumerate}[label=(\alph*),leftmargin=1.1cm]
  \item Calcule a impureza do pai e a impureza ponderada dos filhos, pelas três
        medidas, e obtenha as três reduções.
  \item O corte é bom? (Olhe o filho da direita.) Qual das três medidas
        \emph{não percebe} isso, e qual é exatamente a redução que ela indica?
  \item Que consequência isso tem para o algoritmo que \textbf{cresce} a árvore?
  \item Ainda assim, o erro de classificação é a medida recomendada para
        \textbf{podar}. Por quê?
\end{enumerate}
\end{exercicio}

\begin{solucao}
\textbf{(a)} Proporções: pai $p=0{,}75$, esquerda $p=0{,}50$, direita $p=1{,}00$.
Os pesos são $200/400 = 1/2$ para cada filho.
\begin{center}\small
\renewcommand{\arraystretch}{1.4}
\begin{tabular}{@{}lccc@{}}
\toprule
medida & pai & filhos (ponderado) & redução \\
\midrule
erro       & $\min(0{,}75;0{,}25)=0{,}2500$
           & $\tfrac12(0{,}50)+\tfrac12(0)=0{,}2500$
           & $\mathbf{0{,}0000}$ \\
Gini       & $2(0{,}75)(0{,}25)=0{,}3750$
           & $\tfrac12(0{,}50)+\tfrac12(0)=0{,}2500$
           & $\mathbf{0{,}1250}$ \\
entropia   & $0{,}8113$
           & $\tfrac12(1{,}0000)+\tfrac12(0)=0{,}5000$
           & $\mathbf{0{,}3113}$ \\
\bottomrule
\end{tabular}
\end{center}

\smallskip
\textbf{(b)} O corte é \textbf{claramente bom}: o filho da direita é
\emph{puro} --- $200$ observações, todas da classe~1. Separar 200 observações com
certeza absoluta da classe é exatamente o que uma árvore deveria querer fazer.

A medida que não percebe é o \textbf{erro de classificação}, e a redução que ela
indica é \textbf{exatamente zero}. Não é ``pequena'': é zero, na precisão de
máquina. A razão é que, em ambos os filhos, a classe majoritária continua sendo
a classe~1 --- na esquerda por empate, na direita por unanimidade --- e o erro de
classificação só olha a classe majoritária de cada nó. Ele conta $100$ erros no
pai e $100+0=100$ erros nos filhos.

\smallskip
\textbf{(c)} Se a árvore usasse o erro de classificação como critério de corte,
ela avaliaria este corte como \textbf{inútil} e possivelmente escolheria outro,
ou pararia de crescer (se todos os cortes candidatos dessem redução zero, o que é
comum). Um filho puro seria descartado.

É por isso que o \texttt{DecisionTreeClassifier} do \texttt{scikit-learn} oferece
apenas \texttt{criterion="gini"} e \texttt{"entropy"} (ou \texttt{"log\_loss"})
para crescer --- \textbf{o erro de classificação não está entre as opções}. Gini e
entropia são \emph{estritamente} côncavas, e o Exercício~4 mostra que isso garante
redução positiva sempre que os filhos diferirem.

\smallskip
\textbf{(d)} Porque podar e crescer respondem a perguntas diferentes.

Crescer é uma busca gulosa: precisa de um critério \textbf{sensível}, capaz de
distinguir um corte bom de um corte medíocre mesmo quando nenhum dos dois muda a
classe majoritária. Um critério que devolve zero com frequência não orienta a
busca.

Podar é uma decisão sobre o \textbf{desempenho final} da árvore, e o desempenho
final é medido pelo erro que ela vai cometer --- que é, literalmente, o erro de
classificação. Usar Gini para podar otimizaria uma quantidade que não é a que
interessa ao usuário.

Em uma frase: \emph{Gini e entropia são bons guias, o erro de classificação é o
bom juiz.}
\end{solucao}

%------------------------------------------------------------------------------
\begin{exercicio}[KNN nos extremos, em classificação]
\begin{enumerate}[label=(\alph*),leftmargin=1.1cm]
  \item Qual é o erro de \emph{treino} do KNN com $k=1$, num conjunto sem pontos
        repetidos? E com $k=n$ (supondo $n$ ímpar)?
  \item Descreva a fronteira de decisão nos dois extremos, em termos de viés e
        variância.
  \item Um resultado clássico (Cover e Hart, 1967) diz que, quando
        $n\to\infty$, o erro do classificador de 1 vizinho mais próximo,
        $\risco_{1\text{NN}}$, satisfaz
        \[
          \risco^\ast \;\le\; \risco_{1\text{NN}} \;\le\; 2\,\risco^\ast ,
        \]
        onde $\risco^\ast$ é o erro de Bayes. Se $\risco^\ast=0{,}10$, em que
        intervalo está o erro assintótico do 1-NN? Comente: isso é uma boa ou uma
        má notícia sobre o 1-NN?
  \item Por que, na prática, o KNN com $k$ escolhido por validação cruzada
        costuma bater o 1-NN, se o teorema já garante um fator 2?
\end{enumerate}
\end{exercicio}

\begin{solucao}
\textbf{(a)} Com $k=1$ o erro de treino é \textbf{exatamente zero}: o vizinho
mais próximo de $\X_i$ é o próprio $\X_i$, a distância $0$, então
$\ghat(\X_i)=Y_i$ sempre. Com $k=n$, todos os pontos são vizinhos de todos e o
classificador devolve sempre a classe majoritária do conjunto inteiro; o erro de
treino é a proporção da classe \emph{minoritária}.

\smallskip
\textbf{(b)} Com $k=1$ a fronteira é o contorno do diagrama de Voronoi entre
pontos de classes diferentes: extremamente recortada, com ilhas em torno de
observações isoladas. \textbf{Viés baixo, variância altíssima} --- trocar uma
observação muda a fronteira na vizinhança dela.

Com $k=n$ a fronteira \emph{não existe}: o classificador é constante. \textbf{Viés
máximo, variância zero}. É o análogo em classificação do preditor constante do
Exercício~2 da Lista Teórica~01.

\smallskip
\textbf{(c)} Com $\risco^\ast = 0{,}10$, o erro assintótico do 1-NN está em
$[0{,}10;\ 0{,}20]$.

É uma notícia \textbf{surpreendentemente boa}, e vale apreciar por quê: o 1-NN é
o classificador mais ingênuo que se pode imaginar --- ele não estima nada, não
tem parâmetros, não faz suposição nenhuma sobre a distribuição --- e mesmo assim,
com dados suficientes, nunca erra mais que o \emph{dobro} do classificador ótimo,
que conhece as densidades. Não existe resultado análogo para a regressão linear
ou para o LDA: se as suposições deles falharem, o erro pode ser arbitrariamente
ruim.

A má notícia está na palavra \emph{assintótico}. A Aula~05 mostrou quanto dado é
preciso para chegar perto do regime assintótico em dimensão alta: em $d=10$,
cem milhões de observações para o mesmo risco que cem dariam em $d=1$. A garantia
é real e o $n$ que a torna operante, não.

\smallskip
\textbf{(d)} Por três razões que se somam.

Primeiro, o teorema é \textbf{assintótico}: com $n$ finito o 1-NN está bem acima
do seu próprio limite, porque a variância domina.

Segundo, o fator $2$ é uma \textbf{cota}, e cotas costumam ser folgadas. Um $k$
maior tipicamente chega bem mais perto de $\risco^\ast$ do que de $2\risco^\ast$
--- a média sobre $k$ vizinhos reduz a variância por um fator próximo de $k$, ao
preço de um viés que cresce devagar enquanto os vizinhos ainda estiverem próximos.

Terceiro, e mais importante na prática: o $k$ ótimo depende de $n$, de $d$ e do
nível de ruído, e ninguém sabe qual é. A validação cruzada o encontra sem precisar
conhecer nada disso --- e, como a Lista prática~03 mediu, a curva de erro contra
$k$ costuma ser \emph{plana} numa faixa larga, o que significa que basta acertar
a ordem de grandeza.
\end{solucao}

%------------------------------------------------------------------------------
\begin{exercicio}[o que cada classificador supõe]
Complete a tabela abaixo para os classificadores vistos no curso.
\begin{center}\small
\renewcommand{\arraystretch}{1.35}
\begin{tabular}{@{}lp{4.3cm}cc@{}}
\toprule
classificador & suposição principal & forma da fronteira & gera./discr. \\
\midrule
regressão logística & & & \\
LDA                 & & & \\
QDA                 & & & \\
naive Bayes gaussiano & & & \\
KNN                 & & & \\
árvore              & & & \\
SVM linear          & & & \\
SVM com RBF         & & & \\
\bottomrule
\end{tabular}
\end{center}
\begin{enumerate}[label=(\alph*),leftmargin=1.1cm]
  \item Preencha as três colunas.
  \item Quais desses métodos produzem estimativas de $\Prob(Y=1\mid\X=\x)$
        diretamente? Quais delas você usaria numa conta de custo esperado, à luz
        da Aula~09?
  \item Dois deles são \emph{invariantes} a transformações monótonas de cada
        covariável separadamente (trocar $x_j$ por $\log x_j$, por exemplo). Quais
        são, e por quê? Que consequência prática isso tem?
\end{enumerate}
\end{exercicio}

\begin{solucao}
\textbf{(a)}
\begin{center}\footnotesize
\renewcommand{\arraystretch}{1.35}
\begin{tabular}{@{}lp{5.2cm}ll@{}}
\toprule
classificador & suposição principal & fronteira & tipo \\
\midrule
logística  & $\log\frac{p}{1-p}$ é linear em $\x$; nada sobre a distribuição de $\X$ & linear & discriminativo \\
LDA        & $\X\mid Y=k$ gaussiana, covariância \emph{comum} & linear & generativo \\
QDA        & $\X\mid Y=k$ gaussiana, covariâncias distintas & quadrática & generativo \\
naive Bayes & $\X\mid Y=k$ gaussiana com covariância \emph{diagonal} & quadrática & generativo \\
KNN        & $r$ varia devagar (Lipschitz); a métrica de distância faz sentido & arbitrária, recortada & discriminativo \\
árvore     & $r$ é aproximadamente constante em blocos \emph{alinhados aos eixos} & escada ortogonal & discriminativo \\
SVM linear & fronteira linear; só os pontos da margem importam & linear & discriminativo \\
SVM com RBF & fronteira suave; similaridade decai com a distância euclidiana & suave, arbitrária & discriminativo \\
\bottomrule
\end{tabular}
\end{center}

\smallskip
\textbf{(b)} Dão probabilidades \textbf{diretamente}: logística, LDA, QDA, naive
Bayes e KNN (a proporção entre os $k$ vizinhos). A árvore dá a proporção da folha.
A \textbf{SVM não dá}: ela devolve uma distância com sinal, não uma
probabilidade, e o \texttt{predict\_proba} do \texttt{scikit-learn} só existe com
\texttt{probability=True}, que ajusta uma calibração de Platt \emph{por cima} do
modelo, com validação cruzada interna.

Numa conta de custo esperado (Exercício~2 da Lista Teórica~09) o que importa não
é ter um número entre 0 e 1, é ele ser \textbf{calibrado}. A logística é a mais
confiável nesse quesito, por construção --- ela maximiza a verossimilhança de
Bernoulli, que é uma regra de pontuação própria. O naive Bayes é o pior, pelo
motivo medido na Lista prática~09. KNN e árvores dão probabilidades grosseiras,
com granularidade $1/k$ e $1/n_{\text{folha}}$ respectivamente. Para qualquer um
deles, vale conferir com uma curva de calibração antes de usar o número numa
decisão de custo.

\smallskip
\textbf{(c)} São a \textbf{árvore} e (com uma ressalva) os \emph{ensembles} de
árvores.

A razão é que a árvore só usa as covariáveis por meio de comparações
$x_j < c$. Se $g$ é estritamente crescente, então $x_j < c \iff g(x_j) < g(c)$:
o mesmo conjunto de partições continua disponível, e o corte ótimo é o mesmo com
o limiar transformado. A árvore ajustada é \emph{idêntica}, ponto a ponto.

A consequência prática é grande: \textbf{árvores não precisam de padronização},
nem de correção de assimetria, nem de tratamento de escala. Compare com o LDA (a
matriz de covariância muda), com a SVM com RBF (a distância euclidiana muda) ou
com o KNN (idem) --- todos precisam. É uma das razões de florestas e
\emph{boosting} serem tão populares em dados tabulares heterogêneos: eles
dispensam metade do \texttt{Pipeline} da Aula~07.

(A ressalva: a invariância vale para transformações de \emph{uma} covariável de
cada vez. Rotacionar o espaço --- misturar $x_1$ e $x_2$ --- muda tudo, porque os
cortes são obrigatoriamente alinhados aos eixos. É por isso que árvores vão mal
em fronteiras diagonais, e é exatamente a fragilidade oposta à do LDA.)
\end{solucao}

%------------------------------------------------------------------------------
\begin{exercicio}[concavidade: por que Gini sempre ganha algo]
Seja $\phi:[0,1]\to\R$ uma medida de impureza e considere um nó com proporção
$p$ dividido em dois filhos de pesos $w$ e $1-w$ e proporções $p_E$ e $p_D$, com
\[
  p = w\,p_E + (1-w)\,p_D
\]
(a proporção do pai é sempre a média ponderada das dos filhos --- confira isso no
Exercício~1).
\begin{enumerate}[label=(\alph*),leftmargin=1.1cm]
  \item Mostre que, se $\phi$ é \textbf{côncava}, a redução de impureza
        $\phi(p) - \big[w\,\phi(p_E) + (1-w)\,\phi(p_D)\big]$ é sempre $\ge 0$.
  \item Mostre que, se $\phi$ é \textbf{estritamente} côncava, a redução é
        $>0$ sempre que $p_E \neq p_D$.
  \item Verifique que $\phi(p)=\min(p,1-p)$ é côncava mas \textbf{não}
        estritamente. Use isso para explicar, em uma linha, o resultado do
        Exercício~1.
  \item Mostre que o Gini $2p(1-p)$ é a \textbf{variância} de uma Bernoulli$(p)$.
        Que interpretação isso dá à ``redução de impureza de Gini''?
\end{enumerate}
\end{exercicio}

\begin{solucao}
\textbf{(a)} É exatamente a desigualdade de Jensen. Por definição, $\phi$ côncava
significa
\[
  \phi\big(w\,p_E + (1-w)\,p_D\big) \;\ge\; w\,\phi(p_E) + (1-w)\,\phi(p_D)
\]
para todo $w\in[0,1]$. Como $p = w\,p_E+(1-w)\,p_D$, o lado esquerdo é $\phi(p)$,
e a redução
\[
  \phi(p) - \big[w\,\phi(p_E)+(1-w)\,\phi(p_D)\big] \;\ge\; 0 .
\]
Ou seja: \textbf{cortar nunca aumenta a impureza}, qualquer que seja a medida
côncava. É a garantia mínima que se pede de um critério de corte.

\smallskip
\textbf{(b)} Concavidade estrita é a mesma desigualdade com $>$ no lugar de
$\ge$, desde que $p_E\neq p_D$ e $0<w<1$. Logo a redução é estritamente positiva
nessas condições. Em palavras: \textbf{qualquer corte que separe os filhos em
proporções diferentes é premiado}, mesmo que a classe majoritária não mude.

\smallskip
\textbf{(c)} A função $\min(p,1-p)$ é linear por partes:
\[
  \min(p,1-p) = \begin{cases} p, & 0\le p\le \tfrac12,\\[2pt] 1-p, & \tfrac12\le p\le 1.\end{cases}
\]
É côncava (o gráfico é um ``$\wedge$'' com vértice em $p=1/2$), mas em cada uma
das duas metades ela é \emph{linear} --- e função linear satisfaz Jensen com
igualdade.

No Exercício~1, $p_E = 0{,}50$ e $p_D=1{,}00$ estão \textbf{os dois no ramo
$[\tfrac12, 1]$}, onde $\phi(p)=1-p$ é linear. Então a média ponderada de
$\phi$ é igual a $\phi$ da média ponderada:
\[
  \tfrac12(1-0{,}5)+\tfrac12(1-1{,}0) = 0{,}25 = 1-0{,}75 = \phi(0{,}75),
\]
e a redução é exatamente zero. Não foi coincidência da escolha dos números: é o
que \emph{sempre} acontece quando os dois filhos caem do mesmo lado de $1/2$.

\smallskip
\textbf{(d)} Se $Z\sim\text{Bernoulli}(p)$, então $\E[Z]=p$ e
$\E[Z^2]=\E[Z]=p$ (porque $Z^2=Z$), logo
\[
  \Var(Z) = p - p^2 = p(1-p) = \tfrac12\cdot 2p(1-p) = \tfrac12\,\text{Gini}(p).
\]
O índice de Gini é, portanto, o dobro da variância de uma Bernoulli --- e, a menos
desse fator 2, \emph{é} a variância.

A interpretação é bonita e unifica o curso: a redução de impureza de Gini é
exatamente a \textbf{redução da soma de quadrados dentro dos grupos} que a árvore
de \emph{regressão} da Aula~06 usa, aplicada à variável indicadora $\1{Y=1}$.
Crescer uma árvore de classificação por Gini é a mesma coisa que crescer uma
árvore de regressão sobre o indicador da classe --- o mesmo algoritmo, o mesmo
critério, escrito de outro jeito.
\end{solucao}


\end{document}
